\section{The critical moving-point barrier}
\label{sec:barrier}

We first separate the coefficient of \(\gamma\) from the gamma-free tail.
Fix \(q>1\), put
\[
 f_q(t)=E_1(q^t),
\]
and consider a finitely supported coefficient vector \((c_n)\).  The two
linear statistics
\[
 \omega(c)=\sum_nc_n,
 \qquad
 \mu(c)=\sum_nnc_n
\]
give the exact decomposition
\begin{equation}\label{eq:augmentation-decomposition}
 \sum_nc_n\Ein(q^n)
 =\omega(c)\gamma+\mu(c)\log q+\sum_nc_nf_q(n).
\end{equation}
We call \(\omega=1,\mu=0\) the augmentation-one, log-balanced stratum.
The dyadic level belongs to this stratum.  The next theorem completely solves
its natural cumulative-positive cone.

Direct differentiation gives
\begin{equation}\label{eq:fq-convex}
 f_q'(t)=-(\log q)e^{-q^t}<0,
 \qquad
 f_q''(t)=(\log q)^2q^te^{-q^t}>0.
\end{equation}

\begin{theorem}[Sharp cumulative-positive extremality]
\label{thm:augmentation-extremality}
Let \(M\ge2\), and suppose that real coefficients \(c_0,\ldots,c_M\)
satisfy
\[
 \sum_{n=0}^Mc_n=1,
 \qquad
 \sum_{n=0}^Mnc_n=0.
\]
Put \(C_k=\sum_{n=0}^kc_n\), and assume \(C_k\ge0\) for \(0\le k<M\).
Then the complete sharp range is
\begin{equation}\label{eq:augmentation-range}
 \begin{split}
 M f_q(M-1)-(M-1)f_q(M)
 &\le \sum_{n=0}^Mc_nf_q(n)\\
 &\le f_q(M)+M\{f_q(0)-f_q(1)\}.
 \end{split}
\end{equation}
The lower endpoint is attained uniquely by
\begin{equation}\label{eq:augmentation-minimizer}
 c_{M-1}=M,\qquad c_M=-(M-1),
 \qquad c_0=\cdots=c_{M-2}=0.
\end{equation}
The upper endpoint is attained uniquely by
\(c_0=M,c_1=-M,c_M=1\), with the other coefficients zero.
Consequently the moving level \(Y_{M-1}(q)-\gamma\) is the unique minimum
tail in this cone, and
\[
 \min\sum_nc_nf_q(n)
 \sim \frac{M}{q^{M-1}}e^{-q^{M-1}}.
\]
\end{theorem}

\begin{proof}
Abel summation and the moment normalization give
\begin{align*}
 \sum_{n=0}^Mc_nf_q(n)
 &=f_q(M)+\sum_{k=0}^{M-1}C_k\{f_q(k)-f_q(k+1)\},\\
 \sum_{k=0}^{M-1}C_k&=M.
\end{align*}
By strict convexity, the positive differences
\(f_q(k)-f_q(k+1)\) decrease strictly with \(k\).  The displayed sum is
therefore minimized by putting all mass \(M\) at \(k=M-1\), and maximized
by putting it at \(k=0\).  Unwinding the partial sums gives the two unique
coefficient vectors.  The feasible \(C\)-vectors form the simplex
\(\{C_k\ge0:\sum C_k=M\}\); its continuous linear image is the whole
interval between these two endpoints, which proves the asserted complete
range.  The asymptotic is \eqref{eq:E1-asymptotic}.
\end{proof}

The same doubly exponential scale survives when the two levels use
multiplicatively independent bases.

\begin{proposition}[Two-base moving-level obstruction]
\label{prop:two-base-obstruction}
For \(q>1\), write
\[
 Y_N(q)=(N+1)\Ein(q^N)-N\Ein(q^{N+1})
       =\gamma+\varepsilon_N(q).
\]
Let
\[
 M_N=\left\lfloor\frac{N\log2}{\log3}\right\rfloor+1,
 \qquad W_N=Y_N(2)-Y_{M_N}(3).
\]
Then \(W_N>0\) for all sufficiently large \(N\), and
\begin{equation}\label{eq:two-base-asymptotic}
 \boxed{\displaystyle
 -\log W_N
 =2^N+N\log2-\log(N+1)+O(1).}
\end{equation}
In particular \(-\log W_N=2^N+O(N)\).  Thus changing from one
multiplicative orbit to two independent ones does not remove the
moving-point scale of the fixed-form obstruction.
\end{proposition}

\begin{proof}
Put \(x=2^N\) and \(y=3^{M_N}\).  Since
\(N\log2/\log3\) is not an integer, \(y>x\); both are integers, so
\(y\ge x+1\).  The standard bounds
\[
 \frac{e^{-u}}{u+1}<E_1(u)<\frac{e^{-u}}u
\]
show, uniformly here, that
\[
 \varepsilon_N(2)
 =\frac{N+1}{x}e^{-x}\{1+O(x^{-1})\}.
\]
Moreover \(M_N/N\to\log2/\log3<1\), and the same bounds give
\[
 \frac{\varepsilon_{M_N}(3)}{\varepsilon_N(2)}
 \le
 \left(\frac{\log2}{\log3}+o(1)\right)e^{-1}<1.
\]
Hence \(W_N\) is positive and differs from \(\varepsilon_N(2)\) by a
factor bounded above and below away from zero.  Taking logarithms yields
\eqref{eq:two-base-asymptotic}.
\end{proof}

Outside this cone, augmentation one alone has no sign content.

\begin{proposition}[Three-point flexibility]
\label{prop:three-point-flexibility}
Let \(0<a<b<c\) be positive integers.  Put
\[
 c^{(0)}=\left(\frac b{b-a},-\frac a{b-a},0\right),
 \qquad
 v=(c-b,a-c,b-a).
\]
For \(t\in\R\), the vector \(c(t)=c^{(0)}+tv\), supported at \(a,b,c\),
is augmentation one and log balanced.  Its tail is \(A+tD\), where
\[
 A>0,
 \qquad
 D=(b-a)(c-b)(c-a)f_q[a,b,c]>0.
\]
Hence rational admissible coefficients give a dense set of tail values in
\(\R\); in particular both signs and arbitrarily small nonzero values occur.
For example, at \(q=2\) and exponents \(1,2,3\),
\[
 -E_1(2)+5E_1(4)-3E_1(8)<0.
\]
\end{proposition}

\begin{proof}
The two normalizations are immediate.  The slope \(D\) is a positive second
divided difference by \eqref{eq:fq-convex}.  Since \(A+D\Q\) is dense in
\(\R\), all qualitative assertions follow.  For the explicit example use
\[
 E_1(x)>\frac{e^{-x}}{x+1},
 \qquad
 E_1(x)<\frac{e^{-x}}x:
\]
the displayed value is less than
\(-e^{-2}/3+5e^{-4}/4<0\).
\end{proof}

The next estimate identifies the coefficient height at which such
cancellation can begin.  Write a rational coefficient vector primitively as
\(c_j=a_j/D\), with \(D\ge1\), and put
\(H(c)=\max(D,|a_1|,\ldots,|a_r|)\).

\begin{proposition}[First-tail dominance]
\label{prop:first-tail-dominance}
Let \(n_1<\cdots<n_r\) be integers with \(a_1\ne0\).  If
\begin{equation}\label{eq:height-subcritical}
 (r-1)H(c)\frac{1+q^{-n_1}}q
 e^{-(q-1)q^{n_1}}\le\frac12,
\end{equation}
then the tail has the sign of \(a_1\) and
\begin{equation}\label{eq:first-tail-lower}
 \boxed{\displaystyle
 \left|\sum_{j=1}^rc_jE_1(q^{n_j})\right|
 \ge\frac{e^{-q^{n_1}}}{2H(c)(q^{n_1}+1)}.}
\end{equation}
Thus, for fixed \(q,r\), cancellation of the first active tail requires
\[
 \log H(c)\ge(q-1)q^{n_1}+O(1).
\]
\end{proposition}

\begin{proof}
The reverse triangle inequality gives
\[
 \left|\sum_jc_jE_1(q^{n_j})\right|
 \ge\frac{E_1(q^{n_1})-(r-1)H(c)E_1(q^{n_1+1})}{H(c)}.
\]
The elementary bounds used in the preceding proof imply
\[
 \frac{E_1(q^{n+1})}{E_1(q^n)}
 \le\frac{1+q^{-n}}q e^{-(q-1)q^n}.
\]
Now apply \eqref{eq:height-subcritical}.
\end{proof}

For consecutive triples the transition is exact.  Let
\[
 F_{j,N}=E_1(q^{N+j}),
 \quad
 A_N=(N+2)F_{1,N}-(N+1)F_{2,N},
 \quad
 B_N=F_{0,N}-2F_{1,N}+F_{2,N},
\]
and \(\rho_N=B_N/A_N\).  Both \(A_N,B_N\) are positive.

\begin{theorem}[Exact denominator phase transition]
\label{thm:denominator-transition}
Every rational augmentation-one, log-balanced coefficient vector on
\(N,N+1,N+2\) has the form
\[
 c(s)=(s,N+2-2s,s-N-1).
\]
For \(s=a/D\) with \(D\ge1\), its tail \(\varepsilon_N(s)\) satisfies
\begin{equation}\label{eq:denominator-exact}
 \boxed{\displaystyle
 |\varepsilon_N(a/D)|
 =\frac{B_N}{D}\left|a+\frac D{\rho_N}\right|,}
\end{equation}
and hence
\[
 \min_{a\in\Z}|\varepsilon_N(a/D)|
 =\frac{B_N}{D}\left\|\frac D{\rho_N}\right\|.
\]
In particular, \(D\le\rho_N/2\) implies
\(|\varepsilon_N(a/D)|\ge A_N\) for every integer \(a\).  Conversely,
whenever \(\rho_N\ge\tfrac12\) (hence for all sufficiently large \(N\)),
choosing a positive integer \(D\) nearest to \(\rho_N\) and taking
\(a=-1\) gives either an exact zero or
\begin{equation}\label{eq:denominator-upper}
 0<|\varepsilon_N(-1/D)|\le\frac{A_N}{2D}.
\end{equation}
Furthermore, with \(X=q^N\),
\begin{equation}\label{eq:rho-asymptotic}
 \rho_N\sim\frac q{N+2}e^{(q-1)X}.
\end{equation}
Thus the denominator threshold is exactly
\(e^{(q-1)q^N}/N\) up to a constant factor.
\end{theorem}

\begin{proof}
Substitution gives \(\varepsilon_N(s)=A_N+B_Ns\), proving
\eqref{eq:denominator-exact}.  If \(D/\rho_N\le1/2\), its nearest integer
is zero.  If \(|D-\rho_N|\le1/2\), then
\(|1-D/\rho_N|\le1/(2\rho_N)\), which gives
\eqref{eq:denominator-upper}.  Finally \eqref{eq:E1-asymptotic} yields
\[
 B_N\sim E_1(X),
 \qquad
 A_N\sim(N+2)E_1(qX),
\]
and hence \eqref{eq:rho-asymptotic}.
\end{proof}

When \(q\) is algebraic, exact zero can occur for at most one \(N\); by
\cref{thm:uN-defect}, any such zero would already force the algebraic
independence of \(\gamma,\delta\) over \(K\).  Discarding that possible one
index, the nearest-denominator construction gives normalized coefficient
height
\[
 H_N=q e^{(q-1)X}(1+o(1))
\]
and
\begin{equation}\label{eq:critical-height-upper}
 0<|\varepsilon_N|
 \le\frac{(N+2)^2}{2q^2X}
 e^{-(2q-1)X}(1+o(1)).
\end{equation}
Consequently a uniform lower bound, on normalized nonzero rational
augmentation-one triples, of the form
\[
 |\varepsilon|\ge C e^{-\beta X}H^{-\kappa}X^{-\lambda}
\]
requires
\begin{equation}\label{eq:height-frontier}
 \boxed{\beta+\kappa(q-1)\ge2q-1.}
\end{equation}
At equality it also requires \(\lambda\ge1\).  In particular, at the natural
point cost \(\beta=1\), one must have \(\kappa\ge2\).  This exponent refers
to the affine normalization \(\sum c_j=1\); clearing denominators multiplies
the value and changes the corresponding homogeneous height exponent.

The preceding sections construct exact finite spectral data converging very
rapidly to Euler-level traces.  To extract the arithmetic nature of the
limit, one would need a lower bound uniform in a moving algebraic sample
point.  We first show that sublinear point cost is impossible.

Let
\[
 F_1(z)=\Ein(z),\qquad
 F_2(z)=\Ein(2z),\qquad
 F_3(z)=\Ein(4z),
\]
and let \(P_0=X_1-2X_2+X_3\).

\begin{theorem}[A fixed-form critical-scale no-go]
\label{thm:sublinear-no-go}
For \(x\ge1\),
\begin{equation}\label{eq:R-def}
 R(x):=P_0(F_1(x),F_2(x),F_3(x))
 =E_1(x)-2E_1(2x)+E_1(4x)>0,
\end{equation}
and
\begin{equation}\label{eq:R-asymptotic}
 R(x)=\frac{e^{-x}}x
 \left(1+O(x^{-1})+O(e^{-x})\right).
\end{equation}
Consequently, for every \(A>0\) and every \(\vartheta<1\),
\begin{equation}\label{eq:R-sublinear}
 0<R(n)<e^{-An^\vartheta}
\end{equation}
for all sufficiently large integers \(n\).

More sharply, for every \(C_0>0\), a lower bound
\begin{equation}\label{eq:critical-polynomial-template}
 R(n)\ge C_0e^{-\sigma n}n^{-\lambda}
\end{equation}
cannot hold for all large \(n\) if either \(\sigma<1\), or if
\(\sigma=1\) and \(\lambda<1\).

In particular, no lower bound
\begin{equation}\label{eq:uniform-template}
 |P(F(\xi))|\ge
 \exp\!\left[-C(1+|\xi|)^\vartheta
 \Phi(\deg P,\log H(P))\right]
\end{equation}
can hold uniformly for all positive integral \(\xi\), all nonzero values and
all fixed \(E\)-function systems if \(\vartheta<1\), when \(C\) and \(\Phi\)
are independent of \(\xi\).
\end{theorem}

\begin{proof}
The connection formula \eqref{eq:Ein-connection} cancels \(\gamma\),
\(\log x\), and the scale logarithms because
\[
 1-2+1=0,
 \qquad
 \log x-2\log(2x)+\log(4x)=0.
\]
This proves the identity in \eqref{eq:R-def}.  Using
\(E_1(x)=\int_1^\infty e^{-xt}\,dt/t\), its integrand numerator is, with
\(u=e^{-xt}\),
\[
 u-2u^2+u^4=u(1-u)(1-u-u^2)>0,
\]
because \(0<u\le e^{-1}<(\sqrt5-1)/2\).  This proves positivity.
Asymptotic formula \eqref{eq:R-asymptotic} follows from
\eqref{eq:E1-asymptotic}; the terms at \(2x\) and \(4x\) are exponentially
smaller.  Hence
\[
 -\log R(n)=n+\log n+o(1).
\]
Since \(n^\vartheta=o(n)\), \eqref{eq:R-sublinear} follows.  In
\eqref{eq:uniform-template}, \(P_0\) has fixed degree one and fixed height
two, so its right-hand exponent is \(O(n^\vartheta)\), contradicting
\eqref{eq:R-sublinear}.
Finally, \eqref{eq:R-asymptotic} gives
\[
 \frac{R(n)}{C_0e^{-\sigma n}n^{-\lambda}}
 =C_0^{-1}e^{-(1-\sigma)n}n^{\lambda-1}(1+o(1)),
\]
which tends to zero in each of the two stated cases and contradicts
\eqref{eq:critical-polynomial-template}.
\end{proof}

\begin{remark}[Exact scope of the no-go]
The theorem gives the necessary condition \(\vartheta\ge1\) under the
normalization \eqref{eq:uniform-template}, and at the critical exponential
scale it forces at least the polynomial loss \(n^{-1}\).  It does not prove
the existence or sufficiency of such a bound.  If arbitrary point-dependent
constants are allowed, or if the function system itself is changed with the
sample point, the exponent \(\vartheta\) alone is not an invariant.
\end{remark}

\begin{remark}[Augmentation zero versus one]
The fixed form \(P_0=(1,-2,1)\) has \(\omega(P_0)=0\) and zero logarithmic
moment, so it lies in the gamma-free stratum of
\eqref{eq:augmentation-decomposition}.  It disproves bounds that treat all
coefficient directions uniformly; by itself it does not disprove a theorem
restricted to \(\omega=1\).
\Cref{thm:augmentation-extremality,thm:denominator-transition}
describe what survives, and what fails, after that restriction.
\end{remark}

The same fixed second difference is built into the dyadic tower.

\begin{proposition}[Increment identity]
\label{prop:increment-identity}
For every \(N\ge1\),
\begin{equation}\label{eq:YN-increment}
 \boxed{\displaystyle
 Y_N-Y_{N+1}=(N+1)R(2^N).}
\end{equation}
Consequently
\begin{equation}\label{eq:epsilon-tail-R}
 \varepsilon_N=\sum_{k\ge N}(k+1)R(2^k).
\end{equation}
\end{proposition}

\begin{proof}
Expanding \eqref{eq:YN} gives
\[
 Y_N-Y_{N+1}
 =(N+1)\{\Ein(2^N)-2\Ein(2^{N+1})+\Ein(2^{N+2})\},
\]
which is \eqref{eq:YN-increment}.  Telescoping and
\(Y_N\to\gamma\) give \eqref{eq:epsilon-tail-R}.
\end{proof}

Thus the obstruction to a sublinear measure is not merely analogous to the
tower: it is its exact level increment.

We now formulate the remaining arithmetic question with a coefficient that
is fixed across every level.  For \(s\in\C\), define
\begin{equation}\label{eq:LN-def}
 L_N(s)=2sY_N^2+Y_N-2.
\end{equation}

\begin{theorem}[Common-coefficient critical-slope criterion]
\label{thm:critical-slope}
For \(s\in\C\), the limit
\[
 \lambda(s)=\lim_{N\to\infty}
 \frac{-\log|L_N(s)|}{2^N}
\]
exists whenever the logarithm is defined for all sufficiently large \(N\),
and
\begin{equation}\label{eq:lambda-dichotomy}
 \lambda(s)=
 \begin{cases}
 1,&s=S_2,\\
 0,&s\ne S_2.
 \end{cases}
\end{equation}
For \(s=S_2\), more precisely,
\begin{align}
 L_N(S_2)
 &=\frac{4-\gamma}{\gamma}\varepsilon_N
 +2S_2\varepsilon_N^2>0,\label{eq:LN-exact}\\
 -\log L_N(S_2)
 &=2^N+N\log2-\log(N+1)+O(1).
 \label{eq:LN-log-asymptotic}
\end{align}
Hence
\begin{equation}\label{eq:S2-slope-equivalence}
 \boxed{\displaystyle
 S_2\in\Qbar
 \quad\Longleftrightarrow\quad
 \exists s\in\Qbar\text{ with }\lambda(s)=1.}
\end{equation}
If \(s\in\Qbar\), the family \(\{L_N(s):N\ge1\}\) is countably
algebraically independent over \(K\).
\end{theorem}

\begin{proof}
Since \(2S_2\gamma^2+\gamma-2=0\), substituting
\(Y_N=\gamma+\varepsilon_N\) gives \eqref{eq:LN-exact}; combine this with
\eqref{eq:epsilon-properties} to get \eqref{eq:LN-log-asymptotic}.  If
\(s\ne S_2\), then
\[
 L_N(s)\longrightarrow2\gamma^2(s-S_2)\ne0,
\]
so the normalized logarithm tends to zero.  This proves
\eqref{eq:lambda-dichotomy} and \eqref{eq:S2-slope-equivalence}.

Finally, for algebraic \(s\), one has \(s\in K\) and
\[
 2sY_N^2+Y_N-(L_N(s)+2)=0.
\]
Thus \(K(Y_N:N\in I)\) is algebraic over
\(K(L_N(s):N\in I)\) for every finite set \(I\); when \(s=0\), the relation
is linear.  Algebraic independence of the \(Y_N\) therefore transfers to the
\(L_N(s)\).
\end{proof}

There is an equivalent bounded-distortion form that uses only finite-stage
gamma-free data.

\begin{corollary}[Bounded-distortion criterion]
\label{cor:bounded-distortion}
One has
\begin{equation}\label{eq:bounded-distortion-limit}
 \frac{L_N(S_2)}{Y_N-Y_{N+1}}
 \longrightarrow\frac{4-\gamma}{\gamma},
\end{equation}
whereas for \(s\ne S_2\),
\[
 \left|\frac{L_N(s)}{Y_N-Y_{N+1}}\right|\longrightarrow\infty.
\]
Consequently
\begin{equation}\label{eq:bounded-equivalence}
 S_2\in\Qbar
 \quad\Longleftrightarrow\quad
 \exists s\in\Qbar:\quad
 \left\{\frac{L_N(s)}{Y_N-Y_{N+1}}\right\}_{N\gg1}
 \text{ is bounded}.
\end{equation}
\end{corollary}

\begin{proof}
By \cref{prop:increment-identity} and \eqref{eq:R-asymptotic},
\[
 Y_N-Y_{N+1}\sim\frac{N+1}{2^N}e^{-2^N}
 \sim\varepsilon_N.
\]
Now use \eqref{eq:LN-exact}; if \(s\ne S_2\), the numerator tends to a
nonzero constant while the denominator tends to zero.
\end{proof}

\begin{remark}[Why a generic critical measure is not enough]
The fixed-form obstruction has size \(R(x)\sim e^{-x}/x\).  Under the
algebraicity hypothesis \(s=S_2\), the target at \(x=2^N\) has size
\[
 L_N(S_2)\sim C_\gamma(\log x)\frac{e^{-x}}x.
\]
It is therefore larger, by a logarithmic factor, than an unconditional
fixed-form example already permits.  A lower bound constrained only by this
worst-case pointwise scale cannot distinguish the two.  A surviving route
must use additional structure, for example the repetition of the \emph{same}
algebraic coefficient \(s\) over all levels.
\end{remark}

Existing quantitative \(E\)-function theorems cited here are fixed-point
theorems: the constants in the algebraic-independence measure depend on the
chosen nonzero algebraic evaluation point; see
Adamczewski--Faverjon \cite{AdamczewskiFaverjon2025} and
Fischler--Rivoal \cite{FischlerRivoal2024}.  They do not directly provide
the moving uniformity required by \cref{thm:critical-slope}.  This statement
is only about the scope of those theorems, not an impossibility claim for all
future methods.

\section{An exact asymmetric Sondow integral}
\label{sec:sondow}

One of the discarded barrier heuristics attempted to assign a universal
entropy floor to an asymmetric two-variable Sondow integral.  The correct
replacement is an exact finite formula.  For integers \(A\ge0\) and
\(B\ge1\), put
\begin{equation}\label{eq:Sondow-AB}
 \mathcal I_{A,B}=
 \int_0^1\!\!\int_0^1
 \frac{x^A(1-x)^B y^A(1-y)^B}
 {(1-xy)(-\log(xy))}\,dx\,dy.
\end{equation}

\begin{theorem}[Exact asymmetric evaluation]
\label{thm:Sondow-asymmetric}
One has
\begin{align}
 \mathcal I_{A,B}
 ={}&\binom{2B}{B}\gamma\notag\\
 &+2\sum_{j=0}^B\binom Bj^2
 (H_j-H_{B-j})\log((A+j)!)\notag\\
 &-\sum_{j=0}^B\binom Bj^2H_{A+j}.
 \label{eq:Sondow-exact}
\end{align}
In particular,
\begin{equation}\label{eq:Sondow-example}
 \mathcal I_{1,2}=6\gamma+3\log6-\frac{53}{6}.
\end{equation}
\end{theorem}

\begin{proof}
Use Tonelli's theorem,
\[
 -\frac1{\log(xy)}=\int_0^\infty(xy)^u\,du,
 \qquad
 \frac1{1-xy}=\sum_{k\ge0}(xy)^k,
\]
and the beta integral.  With
\[
 R_B(t)=\frac{B!}{t(t+1)\cdots(t+B)},
\]
this gives
\[
 \mathcal I_{A,B}
 =\sum_{m=A+1}^{\infty}\int_m^\infty R_B(t)^2\,dt.
\]
The rational kernel has the partial fraction identity
\begin{equation}\label{eq:Sondow-partial-fraction}
 \left\{\frac{B!}{t(t+1)\cdots(t+B)}\right\}^{\!2}
 =\sum_{j=0}^B\binom Bj^2
 \left\{\frac1{(t+j)^2}
 +\frac{2(H_j-H_{B-j})}{t+j}\right\}.
\end{equation}
Put \(a_j=\binom Bj^2\) and \(h_j=H_j-H_{B-j}\).  Integrating and summing
only up to \(m=M\) gives the finite identity
\begin{align*}
 \mathcal I_{A,B}^{(M)}
 ={}&\sum_{j=0}^Ba_j(H_{M+j}-H_{A+j})\\
 &-2\sum_{j=0}^Ba_jh_j
   \log\frac{(M+j)!}{(A+j)!}.
\end{align*}
The potentially divergent pieces cancel by
\[
 \sum_ja_j=\binom{2B}{B},\qquad
 \sum_ja_jh_j=0,
\]
together with the derivative of Vandermonde's identity
\[
 2\sum_jja_jh_j=\binom{2B}{B}.
\]
Indeed \(H_{M+j}=\log M+\gamma+o(1)\), while Stirling's formula gives
\[
 \log(M+j)!
 =M(\log M-1)+(j+\tfrac12)\log M
  +\tfrac12\log(2\pi)+o(1).
\]
Letting \(M\to\infty\) in the finite identity gives
\[
 \binom{2B}{B}\gamma
 -\sum_ja_jH_{A+j}
 +2\sum_ja_jh_j\log((A+j)!),
\]
which is \eqref{eq:Sondow-exact}.  Direct substitution gives
\eqref{eq:Sondow-example}.
\end{proof}

Formula \eqref{eq:Sondow-exact} replaces, rather than supports, any fixed
numerical entropy floor for this asymmetric family.  The remaining
Diophantine bottleneck is the transverse height of the prime-exponent vector
in the logarithmic kernel.  We make no novelty claim beyond this exact
evaluation without a separate priority comparison with the broader Sondow
literature \cite{Sondow2003}.

\section{The Pr\'evost grid and residual-prime saturation}
\label{sec:prevost-grid}

For \(0\le m\le n\), Pr\'evost's family \cite{Prevost2008} can be written
\begin{align}
 J_{n,m}&=\gamma+L_{n,m}-2H_n,\label{eq:prevost-form}\\
 L_{n,m}&=\sum_{k=0}^n(-1)^{n+k}
 \binom nk\binom{n+k}k\log(n-m+k+1).
 \label{eq:prevost-log}
\end{align}
Let \(d_n=\lcm(1,\ldots,n)\) and \(j_{n,m}=d_nJ_{n,m}\).  If
\(m/n\to\alpha\in(0,1)\), its saddle rate is
\begin{equation}\label{eq:prevost-rate}
 |J_{n,m}|=\exp\{(\log r(\alpha)+o(1))n\},
 \qquad
 r(\alpha)=
 \alpha^\alpha(1-\alpha)^{2-2\alpha}(2-\alpha)^{\alpha-2},
\end{equation}
up to a polynomial factor.  The unique optimum is
\begin{equation}\label{eq:prevost-optimum}
 \alpha_*=1-\frac1{\sqrt2},
 \qquad
 r(\alpha_*)=3-2\sqrt2,
 \qquad
 -\log r(\alpha_*)=\log(3+2\sqrt2).
\end{equation}

We first isolate an exact linear-algebra obstruction that is independent of
the saddle estimate.  Fix a finite set \(I\) of cells and collect the
normalized forms into
\begin{equation}\label{eq:prevost-vector}
 \mathbf j=\gamma\mathbf B-\mathbf A
 +\sum_{\ell\in\mathcal P_I}(\log\ell)\mathbf V_\ell\in\R^I,
\end{equation}
where \(\mathcal P_I\) is the finite set of primes occurring in the
logarithmic coordinates.  The vectors \(\mathbf B,\mathbf A,\mathbf V_\ell\)
are integral.  We write \(B(x)=\langle\mathbf B,x\rangle\), and similarly
for the other coordinate functionals.

\begin{theorem}[Residual-prime Hilbert saturation]
\label{thm:residual-prime-saturation}
Let \(\mathcal S\subset\mathcal P_I\), let \(a/b\in\Q\) be reduced with
\(b>0\), and put
\[
 E_{\mathcal S,a/b}
 =\bigcap_{\ell\notin\mathcal S}\ker V_\ell
  \cap\ker(aB-bA).
\]
If \(P\) is orthogonal projection onto this subspace, then
\begin{equation}\label{eq:residual-projection}
 \boxed{\displaystyle
 bP\mathbf j=(b\gamma-a)P\mathbf B
 +b\sum_{p\in\mathcal S}(\log p)P\mathbf V_p.}
\end{equation}

Let \(c=(c_p)_{p\in\mathcal S}\in\Z^{\mathcal S}\) be primitive and
nonzero.  Choose \(h_p\in\Z\) with \(\sum h_pc_p=1\), set
\[
 \mathbf T=\sum_{p\in\mathcal S}h_p\mathbf V_p,
 \qquad
 \alpha_c=\sum_{p\in\mathcal S}c_p\log p,
\]
write \(T(x)=\langle\mathbf T,x\rangle\), and define
\[
 E_{c,a/b}=E_{\mathcal S,a/b}
 \cap\bigcap_{p\in\mathcal S}\ker(V_p-c_pT).
\]
Denote its orthogonal projection by \(P_c\).
If \(\gamma=a/b\), then
\begin{equation}\label{eq:residual-direction}
 \boxed{P_c\mathbf j=\alpha_cP_c\mathbf T.}
\end{equation}
If \(\Lambda=E_{c,a/b}\cap\Z^I\) and \(T(\Lambda)=g\Z\), \(g>0\), then
the real quotient minimum
\[
 \lambda_\R=\inf\{\|x\|:x\in E_{c,a/b},\ T(x)=g\}
\]
satisfies
\begin{equation}\label{eq:residual-saturation}
 \boxed{\displaystyle
 \lambda_\R=\frac g{\|P_c\mathbf T\|},
 \qquad
 \lambda_\R\|P_c\mathbf j\|=g|\alpha_c|.}
\end{equation}
\end{theorem}

\begin{proof}
Project \eqref{eq:prevost-vector}.  All primes outside \(\mathcal S\)
vanish and \(aP\mathbf B=bP\mathbf A\), proving
\eqref{eq:residual-projection}.  The fixed-direction equations give
\(P_c\mathbf V_p=c_pP_c\mathbf T\), and
\eqref{eq:residual-direction} follows under \(\gamma=a/b\).
Finally, the restriction of \(T\) to \(E_{c,a/b}\) is represented by
\(P_c\mathbf T\).  Cauchy--Schwarz is sharp in its direction and gives
the first identity in \eqref{eq:residual-saturation}; the second follows
from \eqref{eq:residual-direction}.
\end{proof}

\begin{corollary}[Phase-blind quotient barrier]
\label{cor:phase-blind-quotient}
Under the hypotheses of the second part of
\cref{thm:residual-prime-saturation}, and under \(\gamma=a/b\),
\[
 \lambda_\R\|\mathbf j\|\ge g|\alpha_c|.
\]
Thus a strict inequality in the opposite direction cannot follow from this
real Euclidean quotient identity and Cauchy--Schwarz alone.  Integral
divisibility or genuinely lattice-theoretic information is not excluded.
Any strict reverse inequality would already imply \(\gamma\ne a/b\).
\end{corollary}

\begin{remark}[The one-prime shear]
Let \(P\) now denote orthogonal projection onto
\(\bigcap_{\ell\ne2}\ker V_\ell\), with no candidate rationality
hyperplane imposed, and suppose \(x=P\mathbf V_2\ne0\).  Put
\[
 z=P(\gamma\mathbf B-\mathbf A),
 \qquad \varepsilon=\|P\mathbf j\|.
\]
Then \(P\mathbf j=z+(\log2)x\), and hence
\begin{equation}\label{eq:log2-shear}
 \frac{\det\operatorname{Gram}(x,z)}{\|x\|^2}
 =\|P_{x^\perp}P\mathbf j\|^2\le\varepsilon^2,
 \qquad
 \frac{\langle x,z\rangle}{\|x\|^2}
 =-\log2+O\!\left(\frac\varepsilon{\|x\|}\right).
\end{equation}
Thus \(-\log2\) is the longitudinal shear coefficient, while
\(P_{x^\perp}P\mathbf j\) is the narrow transverse residual.  For an
external proxy \(\widetilde\gamma=a/b\ne\gamma\),
\[
 P(\widetilde\gamma\mathbf B-\mathbf A)
 =z+(\widetilde\gamma-\gamma)P\mathbf B.
\]
More explicitly, if an integral coefficient vector \(c\) kills the
odd-prime labels and has residual label \(k=V_2(c)\), then
\[
 (aB(c)-bA(c))+bk\log2
 =b\langle\mathbf j,c\rangle+(a-b\gamma)B(c).
\]
Thus a proxy form is not controlled by the analytic norm of \(P\mathbf j\)
unless \(a/b=\gamma\); finite-window diagnostics based on an external proxy
cannot be used as irrationality evidence.
\end{remark}

\subsection{The fixed-place Ridout deficit}

For a fixed finite set \(S\) of rational primes and \(0\ne r\in\Z\), write
\[
 r_S=\prod_{p\in S}p^{v_p(r)},
 \qquad r_{S^c}=|r|/r_S.
\]

\begin{proposition}[Fixed-\(S\) Ridout budget]
\label{prop:ridout-budget}
Let \(\xi\) be a real algebraic irrational number.  For every
\(\epsilon>0\), all but finitely many reduced fractions \(a/b\) with
\(b>0\), \(a\ne0\), and \(|\xi-a/b|<1\) satisfy
\begin{equation}\label{eq:ridout-linear-form}
 |b\xi-a|\ge\frac{a_S}{b_{S^c}b^\epsilon}.
\end{equation}
Consequently, suppose along a sequence that
\[
 \log b_N=hN+o(N),
 \quad
 \log(a_N)_S+\log(b_N)_S=\sigma N+o(N),
 \quad
 |b_N\xi-a_N|=e^{-(\beta+o(1))N}.
\]
Assume here that \(h>0\), that \(b_N\to\infty\), and that the reduced
fractions \(a_N/b_N\) take infinitely many distinct values.
If \(\beta>h-\sigma\), then \(\xi\) is not algebraic irrational.  If
\(\beta\le h-\sigma\), these rates alone do not enter the range excluded by
Ridout's theorem.
\end{proposition}

\begin{proof}
Ridout's theorem \cite{Ridout1958} says that
\[
 \prod_{p\in S}|ab|_p
 \min\{1,|\xi-a/b|\}<b^{-2-\epsilon}
\]
has only finitely many reduced solutions.  For a good approximant the left
side is \(|\xi-a/b|/(a_Sb_S)\).  Multiplying the complementary inequality
by \(b\) gives \eqref{eq:ridout-linear-form}; in the homogeneous-height
form of the theorem, \(|\xi-a/b|<1\) makes
\(\max(|a|,b)\ll_\xi b\), and applying the theorem with
\(\epsilon/2\) absorbs the fixed comparison constant for large \(b\).
Taking logarithms and then letting \(\epsilon\downarrow0\) proves the rate
assertion.
\end{proof}

\begin{theorem}[The Pr\'evost saddle--lcm ledger]
\label{thm:prevost-ridout-deficit}
For any integer sequence \(m_n\) with \(m_n/n\to\alpha_*\),
\[
 |d_nJ_{n,m_n}|=
 \exp\{-(\log(3+2\sqrt2)-1+o(1))n\}.
\]
For every fixed finite set \(S\),
\[
 \log(d_n)_S=|S|\log n+O_S(1)=o(n),
 \qquad
 \log d_n=n+o(n).
\]
Thus, in the favorable height-one ledger in which the lcm is the only
denominator cost and a fixed \(S\) contributes no additional exponential
divisibility, the analytic exponent falls short of the Ridout threshold by
\begin{equation}\label{eq:ridout-deficit}
 \boxed{\displaystyle
 \Delta_R=2-\log(3+2\sqrt2)
 =0.2372528259\ldots.}
\end{equation}
This is an exact saddle--denominator deficit, not yet a rational
approximation to \(\gamma\), because logarithmic coordinates remain in
\(J_{n,m_n}\).
\end{theorem}

\begin{proof}
Combine \eqref{eq:prevost-rate} with the prime-number theorem form
\(\log d_n=\psi(n)=n+o(n)\).  Differentiating \(\log r\) gives
\[
 (\log r)'(\alpha)
 =\log\frac{\alpha(2-\alpha)}{(1-\alpha)^2},
\]
which proves the unique optimum \eqref{eq:prevost-optimum}.  Finally
\(v_p(d_n)=\lfloor\log n/\log p\rfloor\), so every fixed \(S\) contributes
only \(o(n)\).  Comparing the resulting exponent
\(\log(3+2\sqrt2)-1\) with the height-one threshold (1) gives
\eqref{eq:ridout-deficit}.
\end{proof}

\begin{corollary}[Conditional rational-extraction budget]
Suppose an elimination of the logarithmic coordinates produces infinitely
many distinct reduced fractions \(a_n/b_n\) such that
\[
 \log b_n=n+o(n),\qquad
 \log(a_n)_S+\log(b_n)_S=o(n),
\]
and preserves the saddle rate
\[
 |b_n\gamma-a_n|
 =\exp\{-(\log(3+2\sqrt2)-1+o(1))n\}.
\]
Then \cref{prop:ridout-budget} does not exclude algebraic irrationality of
\(\gamma\).  Within this height-one template, crossing the Ridout threshold
requires an additional fixed-\(S\) valuation density greater than
\(\Delta_R\), an archimedean gain of the same size, or a corresponding
reduction in height.
\end{corollary}

\begin{remark}[Exact scope]
This is a budget theorem, not an exclusion of every linear combination of
Pr\'evost forms.  Further cancellation could increase the archimedean
exponent, and a new automatic divisibility theorem could increase the
\(S\)-part.  Neither input is presently available.  Enlarging a fixed finite
set \(S\) alone cannot help because \(\log(d_n)_S=o(n)\).
\end{remark}

\begin{proposition}[Bounded-degree logarithmic transfer]
\label{prop:bounded-degree-transfer}
After cancelling all odd-prime coordinates, write a resulting form as
\[
 J=B\gamma-A+k\log2,
 \qquad A,B,k\in\Z,
\]
and assume \(k\ne0\).  If \(\gamma\) is algebraic of degree \(d\), then
\[
 \beta=\frac{A-B\gamma}{k}
\]
has degree at most \(d\), naïve height at most
\(C_\gamma\max(1,|A|,|B|,|k|)^d\), and
\[
 |\log2-\beta|=\frac{|J|}{|k|}.
\]
Thus any such attack is bounded by the degree-\(d\) algebraic approximation
exponent of \(\log2\), with an additional factor \(d\) in the height
exponent.
\end{proposition}

\begin{proof}
If \(f\) is the minimal polynomial of \(\gamma\) and \(B\ne0\), then
\(\beta\) is a zero of
\[
 B^df\!\left(\frac{A-kX}{B}\right),
\]
and the affine transformation shows that its degree is exactly \(d\).
The displayed polynomial has coefficient height
\(O_\gamma(\max(1,|A|,|B|,|k|)^d)\); the standard factor-height bound
gives the same estimate, up to a constant depending on \(\gamma\), for the
minimal polynomial of \(\beta\).  The case \(B=0\) is rational, and the
approximation identity is immediate.
\end{proof}

\begin{remark}
The available bounds are
\[
 \mu(\log2)\le3.574553902525\ldots,
 \qquad
 \mu_2(\log2)\le12.841618132152\ldots;
\]
see \cite{Marcovecchio2009,Marcovecchio2025}.  Thus a rational branch must
beat the first exponent, while the degree-two height transfer makes
\(2\mu_2(\log2)=25.683236264304\ldots\) the corresponding coefficient-height
target.  These are target exponents once the growth of
\(\max(1,|A|,|B|,|k|)\) is known; no comparison with a projected Pr\'evost
family is claimed here without a separate height lemma.
\end{remark}

\section{Fixed-prime Rodrigues twists and universal saddle cancellation}
\label{sec:fixed-prime-rodrigues}

The deficit \eqref{eq:ridout-deficit} suggests inserting a factor
\(p^{cn}\), for one fixed prime \(p\), directly into the Pr\'evost kernel.
The following deformation preserves integrality, degree, and a linear number
of moments, and produces this factor in the endpoint coefficient.  Even if
the whole endpoint factor is favorably credited as nonarchimedean gain, the
real saddle cancels it in the precise sense of
\cref{thm:universal-fixed-p-cancellation}.

Fix a prime \(p\).  For \(n\ge1\), \(1\le\ell\le n\), put
\begin{equation}\label{eq:fixed-p-rodrigues}
 \mathcal P_{n,\ell,p}(x)=
 \frac{(-1)^\ell}{\ell!}\frac{d^\ell}{dx^\ell}
 \left[x^\ell(1-x)^\ell
 \bigl(1+(p-1)x\bigr)^{n-\ell}\right].
\end{equation}
Write \(\mathcal P_{n,\ell,p}(x)=\sum_{k=0}^nc_kx^k\).

\begin{proposition}[Arithmetic and orthogonality]
\label{prop:fixed-p-arithmetic}
One has \(\mathcal P_{n,\ell,p}\in\Z[x]\),
\(\deg\mathcal P_{n,\ell,p}=n\),
\[
 \int_0^1x^j\mathcal P_{n,\ell,p}(x)\,dx=0
 \qquad(0\le j<\ell),
\]
and
\begin{equation}\label{eq:fixed-p-endpoint}
 \boxed{\mathcal P_{n,\ell,p}(1)=p^{n-\ell}.}
\end{equation}
\end{proposition}

\begin{proof}
Expanding before differentiation gives
\[
 c_k=(-1)^\ell\binom{\ell+k}{\ell}
 \sum_{a+b=k}(-1)^a
 \binom\ell a\binom{n-\ell}b(p-1)^b,
\]
which proves integrality and the degree assertion.  Integration by parts
\(\ell\) times proves the moment identities; all boundary terms vanish.
At \(x=1\), only the term in which all derivatives fall on
\((1-x)^\ell\) survives, giving \eqref{eq:fixed-p-endpoint}.
\end{proof}

Put
\[
 W(x)=\frac1{\log x}+\frac1{1-x}
\]
and, for \(0\le m\le\ell\), define
\begin{equation}\label{eq:fixed-p-moment}
 \mathcal J_{n,\ell,m,p}=
 \int_0^1x^{\ell-m}\mathcal P_{n,\ell,p}(x)W(x)\,dx.
\end{equation}

\begin{proposition}[Exact logarithmic linear form]
\label{prop:fixed-p-linear-form}
Let
\[
 Q_{n,\ell,p}=p^{n-\ell},
 \qquad
 A_{n,\ell,p}=\sum_{k=0}^nc_kH_k,
\]
and
\[
 L_{n,\ell,m,p}=
 \sum_{k=0}^nc_k\log(\ell-m+k+1).
\]
Then
\begin{equation}\label{eq:fixed-p-exact-form}
 \boxed{\mathcal J_{n,\ell,m,p}
 =Q_{n,\ell,p}\gamma-A_{n,\ell,p}+L_{n,\ell,m,p},}
\end{equation}
and \(d_nA_{n,\ell,p}\in\Z\).
\end{proposition}

\begin{proof}
Use Euler's integral \(\gamma=\int_0^1W(x)\,dx\) and Frullani's identity
\[
 \int_0^1\frac{x^r-1}{\log x}\,dx=\log(r+1).
\]
The rational part is
\begin{align*}
 &\int_0^1\frac{\mathcal P(1)-x^{\ell-m}\mathcal P(x)}{1-x}\,dx\\
 &\quad=\int_0^1\frac{\mathcal P(1)-\mathcal P(x)}{1-x}\,dx
 +\int_0^1\mathcal P(x)\frac{1-x^{\ell-m}}{1-x}\,dx.
\end{align*}
The second integral vanishes by the moment conditions, and the first is
\(\sum c_kH_k\).  This proves \eqref{eq:fixed-p-exact-form}; its denominator
divides \(d_n\).
\end{proof}

Pr\'evost's positive Markov--Stieltjes representation is
\begin{equation}\label{eq:prevost-markov}
 W(u)=\int_0^1\frac{\omega(t)}{1-(1-u)t}\,dt,
 \qquad
 \omega(t)=\frac1{t\{\log^2(t^{-1}-1)+\pi^2\}}>0.
\end{equation}

\begin{proposition}[Positive double integral]
\label{prop:fixed-p-positive}
One has
\begin{align}
 (-1)^m\mathcal J_{n,\ell,m,p}
 &=\int_0^1\!\int_0^1
 \frac{u^\ell(1-u)^\ell
 (1+(p-1)u)^{n-\ell}
 (1-t)^{\ell-m}t^m}
 {(1-(1-u)t)^{\ell+1}}
 \omega(t)\,dt\,du>0.
 \label{eq:fixed-p-double}
\end{align}
\end{proposition}

\begin{proof}
Insert \eqref{eq:prevost-markov}, use Fubini, and integrate by parts
\(\ell\) times.  The polynomial part is annihilated, while
\[
 \frac{d^\ell}{du^\ell}
 \frac{u^{\ell-m}}{1-(1-u)t}
 =(-1)^m\ell!\frac{(1-t)^{\ell-m}t^m}
 {(1-(1-u)t)^{\ell+1}}.
\]
This gives \eqref{eq:fixed-p-double} and positivity.
\end{proof}

Let \(\theta=\ell/n\), \(\alpha=m/\ell\), and set
\begin{align}
 \phi_\alpha(u,t)&=
 \frac{u(1-u)(1-t)^{1-\alpha}t^\alpha}
 {1-(1-u)t},\label{eq:prevost-saddle-kernel}\\
 \Phi_{\theta,\alpha,p}(u,t)&=
 \phi_\alpha(u,t)^\theta
 (1+(p-1)u)^{1-\theta},\label{eq:fixed-p-saddle}\\
 R_{\theta,\alpha,p}&=-\log\max_{0\le u,t\le1}
 \Phi_{\theta,\alpha,p}(u,t).
\end{align}
We extend \(\phi_\alpha\), and hence
\(\Phi_{\theta,\alpha,p}\), continuously by zero at the boundary corner
\((u,t)=(0,1)\).  Indeed the quotient in
\eqref{eq:prevost-saddle-kernel} tends to zero there for
\(0<\alpha<1\).

For the finite parameters \(\theta_n=\ell/n\),
\(\alpha_n=m/\ell\), the positive double integral is
\[
 (-1)^m\mathcal J_{n,\ell,m,p}
 =\int_0^1\!\int_0^1
 \Phi_{\theta_n,\alpha_n,p}(u,t)^n
 \frac{\omega(t)}{1-(1-u)t}\,dt,du.
\]
The remaining measure has finite mass
\[
 \int_0^1\!\int_0^1
 \frac{\omega(t)}{1-(1-u)t}\,dt,du
 =\int_0^1W(u)\,du=\gamma.
\]
This gives the upper Laplace bound, while a small rectangle around an
interior maximizer gives the matching lower bound.  Consequently,
whenever \(\ell/n\to\theta\in(0,1]\) and
\(m/\ell\to\alpha\in(0,1)\),
\begin{equation}\label{eq:fixed-p-laplace}
 \lim_{n\to\infty}|\mathcal J_{n,\ell,m,p}|^{1/n}
 =e^{-R_{\theta,\alpha,p}}.
\end{equation}

\begin{theorem}[Universal fixed-prime saddle cancellation]
\label{thm:universal-fixed-p-cancellation}
For every prime \(p\), \(0<\theta\le1\), and \(0<\alpha<1\),
\begin{equation}\label{eq:universal-fixed-p-bound}
 \boxed{\displaystyle
 R_{\theta,\alpha,p}+(1-\theta)\log p
 \le\log(3+2\sqrt2).}
\end{equation}
Thus no member of \eqref{eq:fixed-p-rodrigues} improves the original
Pr\'evost saddle after the complete endpoint factor
\(p^{n-\ell}\) is favorably credited as fixed-\(p\) gain.
\end{theorem}

\begin{proof}
Since
\[
 e^{-\{R_{\theta,\alpha,p}+(1-\theta)\log p\}}
 =\max_{u,t}
 \left[
  \phi_\alpha(u,t)^\theta
  \left(\frac{1+(p-1)u}{p}\right)^{1-\theta}
 \right],
\]
it is enough to bound this normalized maximum from below.  Now
For \(0\le u\le1\),
\[
 \frac{1+(p-1)u}{p}=u+\frac{1-u}{p}\ge u.
\]
At \(t=1/2\), the dependence on \(\alpha\) disappears:
\[
 \phi_\alpha\!\left(u,\frac12\right)
 =\frac{u(1-u)}{1+u}.
\]
Consequently the normalized maximum is at least
\[
 \max_{0\le u\le1}
 u\left(\frac{1-u}{1+u}\right)^\theta
 \ge\max_{0\le u\le1}\frac{u(1-u)}{1+u}
 =3-2\sqrt2.
\]
Taking negative logarithms proves \eqref{eq:universal-fixed-p-bound}.
\end{proof}

\begin{corollary}[The fixed-place deficit survives]
\label{cor:fixed-p-deficit}
Even under the favorable assumptions that logarithmic coordinates are
removed at no exponential cost, \(d_n=e^{n+o(n)}\) is the only denominator
cost, and the full \(p^{n-\ell}\) factor survives primitive reduction, the
integerized exponent satisfies
\[
 \bigl(R_{\theta,\alpha,p}-1\bigr)+(1-\theta)\log p
 \le\log(3+2\sqrt2)-1<1.
\]
Thus, even after the entire endpoint factor is credited, this favorable
ledger does not cross the fixed-\(S\) Ridout threshold; its remaining gap is
at least \(\Delta_R\) from \eqref{eq:ridout-deficit}.
\end{corollary}

\begin{proposition}[Balanced factorial ratios have zero fixed-prime slope]
\label{prop:balanced-factorial-fixed-p}
Let
\[
 \mathcal R(n,k)=
 \frac{\prod_{i=1}^rL_i(n,k)!}{\prod_{j=1}^sM_j(n,k)!},
\]
where the integral linear forms are nonnegative and \(O(n)\) on the index
range, and suppose \(\sum_iL_i=\sum_jM_j\).  Then for every fixed prime
\(p\),
\[
 \lvert v_p(\mathcal R(n,k))\rvert=O(\log n)
\]
uniformly wherever the ratio is defined.  In particular its fixed-prime
valuation density is zero.
\end{proposition}

\begin{proof}
Legendre's formula gives a sum over \(p^\nu\).  At each level the linear
parts cancel by balance, leaving a bounded floor error; only
\(O(\log n)\) levels are nonzero.
\end{proof}

\begin{remark}
This covers the standard balanced Ap\'ery, Franel, Domb, and Christol
factorial ratios.  The valuation exponents supplied directly by the
digit-count and \(p\)-Lucas mechanisms of \cite{Delaygue2018} involve at
most \(1+\lfloor\log_p n\rfloor\) base-\(p\) digits and therefore provide
only an \(O(\log n)\) guaranteed gain.  This does not exclude additional
problem-specific cancellation inside a multisum.  If a factorial-balance
defect is \(\delta n+O(1)\), \(\delta\ne0\), Stirling's formula produces a
\(\delta n\log n\) term in the archimedean ledger.  A successful kernel
must therefore use information not already encoded by the balanced
factorial ratio or the positive endpoint weight twist analyzed here.
\end{remark}

\section{The Stokes-framing barrier}
\label{sec:stokes-framing}

The preceding barriers distinguish gamma-free directions from the Euler
extension.  The analogous distinction for an elementary differential
operator annihilating the Euler \(E\)-function is exact.

Put
\[
 \mathscr E(z)=\sum_{n\ge1}\frac{z^n}{n\,n!},
 \qquad F(x)=\mathscr E(-x)=-\Ein(x),
\]
and let \(\partial=\partial_x\).  Consider
\begin{equation}\label{eq:euler-E-operator}
 \mathcal L=(\partial+1)\circ\partial\circ x\circ\partial,
 \qquad
 \mathcal Ly=xy'''+(x+2)y''+y'.
\end{equation}

\begin{theorem}[Euler Stokes-framing barrier]
\label{thm:stokes-framing}
On \(\Omega=\C\setminus(-\infty,0]\), with principal branches,
\[
 \mathcal B_\infty=(1,\log x,E_1(x)),
 \qquad
 \mathcal B_0=(1,F(x),\log x)
\]
are solution bases of \(\mathcal L\), and
\begin{equation}\label{eq:C-gamma}
 \begin{pmatrix}1\\F(x)\\\log x\end{pmatrix}
 =C_\gamma
 \begin{pmatrix}1\\\log x\\E_1(x)\end{pmatrix},
 \qquad
 C_\gamma=
 \begin{pmatrix}
 1&0&0\\
 -\gamma&-1&-1\\
 0&1&0
 \end{pmatrix}.
\end{equation}
The complete upper/lower lateral transition across the negative real axis is
\begin{equation}\label{eq:stokes-transition}
 \begin{pmatrix}1\\\log_+x\\E_{1,+}(x)\end{pmatrix}
 =S
 \begin{pmatrix}1\\\log_-x\\E_{1,-}(x)\end{pmatrix},
 \qquad
 S=
 \begin{pmatrix}
 1&0&0\\
 2\pi i&1&0\\
 -2\pi i&0&1
 \end{pmatrix}.
\end{equation}
In particular \(S\) is independent of \(\gamma\), and the coefficient row
\(c_\gamma=(-\gamma,-1,-1)\) satisfies
\begin{equation}\label{eq:stokes-cocycle-blind}
 \boxed{c_\gamma S=c_\gamma}
\end{equation}
for every complex value of \(\gamma\).

For every \(t\in\Qbar\), \(F+t\) is an \(E\)-function solution of the same
operator.  With \(\mathcal B_0^{(t)}=(1,F+t,\log x)\), its connection
matrix is
\begin{equation}\label{eq:framing-shear}
 C_t=U_tC_\gamma,
 \qquad
 U_t=\begin{pmatrix}1&0&0\\t&1&0\\0&0&1\end{pmatrix}.
\end{equation}
Consequently, under \(\gamma\in\Qbar\), the choice \(t=\gamma\) removes
the Euler entry by an algebraic change of Frobenius framing without changing
the operator or any unframed Stokes conjugacy class.
\end{theorem}

\begin{proof}
Since \(xF'(x)=e^{-x}-1\),
\[
 (\partial+1)\partial(xF')=(\partial+1)(-e^{-x})=0.
\]
Moreover,
\[
 W(1,\log x,E_1(x))=\frac{e^{-x}}{x^2}\ne0,
\]
and \(\det C_\gamma=1\), so both displayed triples are indeed bases.
The connection identity
\[
 F(x)=-\gamma-\log x-E_1(x)
\]
proves \eqref{eq:C-gamma}.  For \(r>0\),
\[
 \log_+(-r)-\log_-(-r)=2\pi i,
 \qquad
 E_{1,+}(-r)-E_{1,-}(-r)=-2\pi i,
\]
which gives \eqref{eq:stokes-transition}.  Matrix multiplication proves
\eqref{eq:stokes-cocycle-blind}.  Adding the constant solution \(t\) adds
\(t\) times the first row of the connection matrix to its second row,
proving \eqref{eq:framing-shear}.
\end{proof}

\begin{corollary}[Exterior blindness]
\label{cor:stokes-exterior-blind}
One has
\[
 \det C_\gamma=1,
 \qquad \det S=1,
 \qquad \operatorname{tr}S=3.
\]
The displayed lateral transition and every invariant of its unframed
conjugacy class are independent of \(\gamma\).  Their determinant, trace,
and exterior-power characteristic data lose the off-diagonal Euler framing.
\end{corollary}

\begin{remark}[Strict factors and framed Stokes constants]
Fischler--Rivoal use ``Stokes constants'' for sector-dependent coefficients
of a distinguished \(E\)-function in a formal basis at infinity; in that
broader sense \(-\gamma\) is such a coefficient.  After the logarithmic
formal monodromy is separated, the strict factor is
\(E_1\mapsto E_1-2\pi i\) and is gamma-free.  Their arithmetic theorem is a
containment theorem for possible coefficients, not a numerical-to-motivic
separation theorem; see \cite{FischlerRivoal2016}.
\end{remark}

\begin{remark}[Class-group averaging does not erase \(\gamma\)]
\label{rem:Chowla-Selberg-correction}
Suppose a Kronecker--Chowla--Selberg family is written in class-group
components as \(A_C=\kappa\gamma+B_C\), with the same Euler coefficient in
each class.  Character projection gives
\[
 \sum_C\overline{\chi(C)}A_C
 =\kappa\gamma\sum_C\overline{\chi(C)}
 +\sum_C\overline{\chi(C)}B_C.
\]
Thus the trivial character, equivalently the class-group average, preserves
the Euler term.  It disappears only in nontrivial character components or
other coefficient-sum-zero directions.  In representation-theoretic
language, \(\gamma\) occupies the trivial isotypic component.  This corrects
the opposite averaging heuristic; it is not a new Chowla--Selberg theorem.
See \cite{GrossCS}.
\end{remark}

\section{The mixed-Gevrey aperture}
\label{sec:mixed-gevrey}

Fischler--Rivoal call a formal Laurent series mixed arithmetic-Gevrey when
its nonnegative part is an \(E\)-function and its negative part is an
arithmetic Gevrey series of order one
\cite{FischlerRivoalMixed2024}.  The Euler connection formula produces a
particularly transparent family in this class.

For \(\lambda\in\Qbar_{>0}\), set
\begin{equation}\label{eq:mixed-dilation-family}
 \widehat D_\lambda(z)=
 \sum_{n\ge0}(-1)^n n!\lambda^{-n-1}z^{-n-1},\qquad
 U_\lambda(z)=e^{\lambda z},\qquad
 W_\lambda(z)=e^{\lambda z}\Ein(\lambda z)-\widehat D_\lambda(z).
\end{equation}

\begin{proposition}[Dilation system and functional rank]
\label{prop:mixed-dilation-rank}
Direction zero is not anti-Stokes for \(\widehat D_\lambda\), and its
Borel sum is
\begin{equation}\label{eq:Dhat-Borel-sum}
 (\widehat D_\lambda)_0(z)=e^{\lambda z}E_1(\lambda z).
\end{equation}
Consequently
\begin{equation}\label{eq:Wlambda-sum}
 (W_\lambda)_0(z)=e^{\lambda z}
 (\gamma+\log\lambda+\log z),
\end{equation}
on the principal positive sector, and
\begin{equation}\label{eq:mixed-dilation-system}
 \binom{U_\lambda}{W_\lambda}'
 =\begin{pmatrix}\lambda&0\\z^{-1}&\lambda\end{pmatrix}
 \binom{U_\lambda}{W_\lambda}.
\end{equation}
For distinct \(\lambda_1,\ldots,\lambda_m\in\Qbar_{>0}\), the \(2m\)
mixed functions \(U_{\lambda_i},W_{\lambda_i}\) are linearly independent
over \(\C(z)\).
\end{proposition}

\begin{proof}
The Laplace integral of the formal negative part is
\[
 \int_0^\infty\frac{e^{-\lambda zt}}{1+t}\,dt
 =e^{\lambda z}E_1(\lambda z),
\]
which proves \eqref{eq:Dhat-Borel-sum}; its Borel singularity lies on the
negative direction.  Subtracting this identity from
\(e^{\lambda z}\Ein(\lambda z)\) and using
\eqref{eq:Ein-connection} proves \eqref{eq:Wlambda-sum}.  Direct
differentiation proves \eqref{eq:mixed-dilation-system}.

Suppose a rational-function linear relation among all the displayed
functions exists, and sum it in direction zero.  Analytic continuation once
around zero adds \(2\pi i e^{\lambda_i z}\) to
\((W_{\lambda_i})_0\) and fixes \(U_{\lambda_i}\).  Subtraction gives a
rational-function relation among exponentials with distinct exponents, so
all coefficients of the \(W\)-terms vanish.  The remaining exponential
relation then forces all \(U\)-coefficients to vanish.
\end{proof}

\begin{proposition}[Unconditional non-divisibility]
\label{prop:mixed-nondivisibility}
For every \(\lambda\in\Qbar_{>0}\) and \(a\in\Qbar\), the formal quotient
\begin{equation}\label{eq:mixed-quotient}
 \frac{W_\lambda-aU_\lambda}{z-1}
\end{equation}
is not a mixed arithmetic-Gevrey function.
\end{proposition}

\begin{proof}
Write the normalized positive part of the numerator as
\[
 F(z)=e^{\lambda z}\{\Ein(\lambda z)-a\}
\]
and its negative part as \(-\widehat D_\lambda(z)\), which has zero constant
term.  If the quotient were mixed, comparison of nonnegative and
nonpositive Laurent parts after multiplication by \(z-1\) would force
\(F(1)\in\Qbar\).  But
\[
 F(1)=e^\lambda\{\Ein(\lambda)-a\},
\]
which is transcendental because \(\Ein(\lambda)\) is algebraically
independent over \(K\supset\Qbar(e^\lambda)\) by
\cref{thm:input-independence}.  This is a contradiction.
\end{proof}

We can now state the exact conditional payoff.  Write
\(\mathbf E,\mathbf D\) for the rings of algebraic-point values of
\(E\)-functions and of directionally summed arithmetic Gevrey-one series,
respectively, in the notation of Fischler--Rivoal.

\begin{theorem}[Conditional simultaneous dilation theorem]
\label{thm:conditional-mixed-dilations}
Assume the intersection conjecture
\begin{equation}\label{eq:ED-intersection}
 \mathbf E\cap\mathbf D=\Qbar.
\end{equation}
For any distinct
\(\lambda_1,\ldots,\lambda_m\in\Qbar_{>0}\), the \(2m\) numbers
\begin{equation}\label{eq:mixed-dilation-values}
 \boxed{\displaystyle
 e^{\lambda_i},\qquad
 e^{\lambda_i}(\gamma+\log\lambda_i)
 \quad(1\le i\le m)}
\end{equation}
are linearly independent over \(\Qbar\).  In particular every
\(\gamma+\log\lambda_i\) is transcendental; the case \(\lambda=1\) gives
the conditional transcendence of \(\gamma\).
\end{theorem}

\begin{proof}
Assume that the values in \eqref{eq:mixed-dilation-values} satisfy a
linear relation over \(\Qbar\).  Using
\(W_\lambda=e^{\lambda z}\Ein(\lambda z)-\widehat D_\lambda\), move all
summed Gevrey terms to the right.  The left side is in \(\mathbf E\), the
right side is in \(\mathbf D\), so their common value is algebraic by
\eqref{eq:ED-intersection}.  The left side has the form
\[
 \sum_i a_i e^{\lambda_i}
 +\sum_i b_i e^{\lambda_i}\Ein(\lambda_i).
\]
The \(\Ein(\lambda_i)\) are algebraically independent over the field \(K\)
containing all \(e^{\lambda_i}\), so algebraicity forces every \(b_i=0\).
The Lindemann--Weierstrass theorem applied to
\(0,\lambda_1,\ldots,\lambda_m\) then forces every \(a_i=0\).  This proves
linear independence.  The final transcendence assertion follows because
an algebraic ratio would make the two values in its block linearly
dependent.
\end{proof}

The dilation construction has an all-order resonant extension.  Write
\(\stirling{n}{r}\) for the unsigned Stirling number of the first kind, and for
\(r\ge1\) put
\begin{equation}\label{eq:all-order-Dhat}
 \widehat D_{\lambda,r}(z)=
 \sum_{n\ge0}(-1)^n\stirling{n+1}{r}
 \lambda^{-n-1}z^{-n-1},
\end{equation}
and
\begin{equation}\label{eq:all-order-mixed-family}
 \mathcal W_{\lambda,0}(z)=e^{\lambda z},\qquad
 \mathcal W_{\lambda,r}(z)
 =e^{\lambda z}\Ein_r(\lambda z)-\widehat D_{\lambda,r}(z)
 \quad(r\ge1).
\end{equation}

\begin{proposition}[All-order mixed jet]
\label{prop:all-order-mixed-jet}
Every \(\mathcal W_{\lambda,r}\) is mixed arithmetic-Gevrey, direction zero
is non-anti-Stokes, and
\begin{equation}\label{eq:all-order-mixed-generating}
 \boxed{\displaystyle
 \sum_{r\ge0}(\mathcal W_{\lambda,r})_0(z)t^r
 =e^{\lambda z}\Gamma(1-t)(\lambda z)^t.}
\end{equation}
They satisfy the triangular differential system
\begin{equation}\label{eq:all-order-mixed-system}
 \mathcal W_{\lambda,0}'=\lambda\mathcal W_{\lambda,0},
 \qquad
 \mathcal W_{\lambda,r}'
 =\lambda\mathcal W_{\lambda,r}+z^{-1}\mathcal W_{\lambda,r-1}
 \quad(r\ge1).
\end{equation}
For distinct positive algebraic \(\lambda_i\), every finite family of the
functions \(\mathcal W_{\lambda_i,r}\) is linearly independent over
\(\C(z)\).
\end{proposition}

\begin{proof}
The normalized Borel transform of \eqref{eq:all-order-Dhat} is
\[
 \frac{(-1)^{r-1}}{\lambda(r-1)!}
 \frac{\log^{r-1}(1+\xi/\lambda)}{1+\xi/\lambda},
\]
a \(G\)-function whose singularities lie on the negative ray.  Hence the
formal series is arithmetic Gevrey of order one and direction zero is
admissible.  The elementary split
\[
 \Gamma(1-t)x^t
 =1+t\sum_{r\ge1}\Ein_r(x)t^{r-1}
   -t x^t\Gamma(-t,x)
\]
together with the large-
\(x\) expansion
\[
 e^x x^t\Gamma(-t,x)
 \sim\sum_{n\ge0}(-1)^n(t+1)_n x^{-n-1}
\]
proves \eqref{eq:all-order-mixed-generating}, since
\([t^{r-1}](t+1)_n=\stirling{n+1}{r}\).  Differentiating the generating
identity and comparing coefficients proves
\eqref{eq:all-order-mixed-system}.

The sectorial sum on the right of
\eqref{eq:all-order-mixed-generating} is \(e^{\lambda z}\) times a
polynomial in \(\log z\) of degree \(r\), with nonzero leading coefficient
\(1/r!\).  Repeated monodromy differences about zero strip off the highest
logarithmic degree.  At each step the linear independence of exponentials
with distinct exponents forces the corresponding rational coefficients to
vanish.  Descending induction on \(r\) proves the last assertion.
\end{proof}

\begin{theorem}[Conditional all-order Gamma-jet independence]
\label{thm:conditional-all-order-gamma}
Assume \eqref{eq:ED-intersection}.  Let
\(\lambda_1,\ldots,\lambda_m\in\Qbar_{>0}\) be distinct and choose
integers \(R_i\ge0\).  Then all numbers
\begin{equation}\label{eq:conditional-all-order-values}
 \boxed{\displaystyle
 e^{\lambda_i}[t^r]\{\Gamma(1-t)\lambda_i^t\},
 \qquad 1\le i\le m,\quad0\le r\le R_i,}
\end{equation}
are linearly independent over \(\Qbar\).  In particular
\begin{equation}\label{eq:conditional-gamma-derivative-LI}
 \boxed{1,\Gamma'(1),\Gamma''(1),\ldots}
\end{equation}
are linearly independent over \(\Qbar\).
\end{theorem}

\begin{proof}
In a putative relation, split every \(r\ge1\) value as the value of the
\(E\)-function \(e^{\lambda z}\Ein_r(\lambda z)\) minus the summed
Gevrey value \((\widehat D_{\lambda,r})_0\), and move all Gevrey terms to
one side.  By \eqref{eq:ED-intersection}, the common value is algebraic.
The other side is a linear form in the jointly algebraically independent
coordinates \(\Ein_r(\lambda_i)\) over \(K\), by
\eqref{eq:resonant-input}; therefore every coefficient with \(r\ge1\)
vanishes.  Lindemann--Weierstrass then removes the remaining \(r=0\)
exponential terms.  Formula \eqref{eq:all-order-mixed-generating} at
\(z=1\) proves \eqref{eq:conditional-all-order-values}.  Taking
\(\lambda=1\) and using
\([t^r]\Gamma(1-t)=(-1)^r\Gamma^{(r)}(1)/r!\) proves
\eqref{eq:conditional-gamma-derivative-LI}.
\end{proof}

\begin{remark}[Priority and scope]
Fischler--Rivoal already prove from \eqref{eq:ED-intersection} the
transcendence of Gamma derivatives at the relevant rational points; in
particular the one-point transcendence of \(\gamma\) is not claimed as new
here.  Their stronger mixed zero-removal Conjecture~3 implies
\eqref{eq:ED-intersection}; alternatively, their conditional linear
specialization theorem and \cref{prop:mixed-dilation-rank} give a second
proof of the displayed rank.  The contribution of
\cref{thm:conditional-mixed-dilations,thm:conditional-all-order-gamma} is
the explicit simultaneous dilation family, its resonant all-order lift, the
use of the companion multipoint theorem, and the full linear rank.  It
does not assert algebraic independence: products of mixed functions need
not be mixed, so the cited specialization theorem is deliberately linear.
The value \((W_1)_0(1)=e\gamma\) is the same Euler extension coordinate
that reappears in the framed motive below, but no implication between the
two conjectural routes is known.
\end{remark}

\section{The Euler--Tate matched-line target}
\label{sec:euler-tate}

The exponential-motive construction supplies the canonical framing absent
from the scalar Stokes conjugacy class.

\begin{theorem}[Euler--Mascheroni motive; Fres\'an--Jossen]
\label{thm:FJ-euler-motive}
There is a rank-two exponential motive \(M(\gamma)\) with adapted de Rham
and Betti bases \((d_0,d_1)\) and \((b_0,b_1)\) for which, with columns
indexed by the de Rham basis and rows by the Betti basis, the comparison
matrix is
\begin{equation}\label{eq:euler-period-matrix}
 P_\gamma=\begin{pmatrix}1&\gamma\\0&2\pi i\end{pmatrix}.
\end{equation}
Equivalently,
\[
 \operatorname{comp}(d_0)=b_0,
 \qquad
 \operatorname{comp}(d_1)=\gamma b_0+2\pi i\,b_1.
\]
It lies in a non-split exact sequence
\begin{equation}\label{eq:euler-motive-extension}
 0\longrightarrow\Q(0)\longrightarrow M(\gamma)
 \longrightarrow\Q(-1)\longrightarrow0,
\end{equation}
satisfies \(\operatorname{Hom}(M(\gamma),\Q(0))=0\), and has motivic
Galois group
\[
 G_{M(\gamma)}=
 \left\{\begin{pmatrix}1&b\\0&a\end{pmatrix}:
 a\in\mathbb G_m,\ b\in\mathbb G_a\right\}.
\]
\end{theorem}

\begin{proof}
This is the Euler--Mascheroni construction and Galois-group computation of
Fres\'an--Jossen \cite[\S12.8]{FresanJossen2024}.
\end{proof}

Its determinant is \(2\pi i\), the same as that of the semisimplification,
and is again blind to the extension coordinate.  The missing input can be
stated for this one object.

\begin{proposition}[Complete matched-line classification]
\label{prop:matched-line-classification}
Every de Rham line mapping isomorphically to the quotient is uniquely of
the form
\[
 W_{\mathrm{dR}}(x)=\Qbar(xd_0+d_1),\qquad x\in\Qbar,
\]
and every Betti line with the same property is uniquely of the form
\[
 W_{\mathrm B}(u)=\Qbar(ub_0+b_1),\qquad u\in\Qbar.
\]
The two lines are matched if and only if
\begin{equation}\label{eq:matched-line-equation}
 \boxed{x+\gamma=2\pi i\,u.}
\end{equation}
Consequently a matched pair exists if and only if
\begin{equation}\label{eq:matched-line-locus}
 \boxed{\gamma\in\Qbar+2\pi i\,\Qbar
              =\Qbar+\pi\Qbar.}
\end{equation}
\end{proposition}

\begin{proof}
The displayed parameterizations are the affine charts of lines transverse
to the respective subobjects.  Comparison sends
\[
 xd_0+d_1\longmapsto(x+\gamma)b_0+2\pi i\,b_1.
\]
This spans \(ub_0+b_1\) exactly when
\eqref{eq:matched-line-equation} holds.  Eliminating \(x,u\in\Qbar\)
gives \eqref{eq:matched-line-locus}; the two fields coincide because
\(2i\in\Qbar^\times\).
\end{proof}

\begin{problem}[Euler--Tate matched-line lifting]
\label{prob:euler-tate-fullness}
After extension to \(\Qbar\), suppose one-dimensional de Rham and Betti
subspaces both map isomorphically onto the corresponding realization of
\(\Q(-1)\), and comparison maps one line to the other.  Prove that the
matched pair is induced by a subobject
\[
 \Q(-1)_{\Qbar}\hookrightarrow M(\gamma)_{\Qbar}.
\]
\end{problem}

\begin{theorem}[Strengthened payoff of matched-line lifting]
\label{thm:matched-line-implies-gamma}
A positive solution of \cref{prob:euler-tate-fullness} implies that
\begin{equation}\label{eq:matched-line-payoff}
 \boxed{\gamma\notin\Qbar+\pi\Qbar.}
\end{equation}
Equivalently, \(1,\pi,\gamma\) are linearly independent over \(\Qbar\).
In particular \(\gamma\), \(\gamma/\pi\), and \(S_2\) are transcendental.
\end{theorem}

\begin{proof}
If \(\gamma\in\Qbar+\pi\Qbar\),
\cref{prop:matched-line-classification} supplies a matched pair.  The
proposed lifting would give a copy of \(\Q(-1)\) mapping nontrivially to the quotient
in \eqref{eq:euler-motive-extension}.  The composite is a nonzero scalar,
so rescaling gives a section and hence a splitting.  Moreover
\(\operatorname{Hom}(-,-)\) commutes here with extension from \(\Q\) to
\(\Qbar\), so the Hom-vanishing in \cref{thm:FJ-euler-motive} remains
zero.  This contradiction proves \eqref{eq:matched-line-payoff}.  Membership
in \(\Qbar+\pi\Qbar\) is exactly a nontrivial \(\Qbar\)-linear relation
among \(1,\pi,\gamma\).  Algebraicity of either \(\gamma\) or
\(\gamma/\pi\) would give such a relation.  Finally, if \(S_2\) were
algebraic, \(2S_2\gamma^2+\gamma-2=0\) would make \(\gamma\) algebraic.
\end{proof}

\begin{remark}[Rational and algebraic versions]
A matched-line lifting statement restricted to \(\Q\)-defined lines would
already exclude the corresponding rational matched locus.  The
\(\Qbar\)-coefficient formulation in \cref{prob:euler-tate-fullness} is the
version that yields the full linear-independence conclusion
\eqref{eq:matched-line-payoff}.
\end{remark}

\begin{remark}[Exact scope of current rapid-decay structures]
Snodgrass constructs cyclotomic rational structures on rapid-decay
cohomology under a hypothesis chosen so that Stokes constants do not enter
\cite[\S1.2]{Snodgrass2026}.  This supplies a realization and period pairing,
but not the numerical-to-motivic fullness asserted in
\cref{prob:euler-tate-fullness}.  The problem is a single-object rank-one
target, far more narrowly scoped than the full exponential period conjecture
but not presently proved.
\end{remark}

\begin{remark}[Three independent apertures]
The common-coefficient problem \cref{prob:common-coefficient} is quantitative
and uses one algebraic coefficient across infinitely many moving levels.
The matched-line problem is categorical and uses one canonical framed
extension.  The mixed-Gevrey route uses the intersection
\(\mathbf E\cap\mathbf D=\Qbar\).  Each would prove the transcendence of
\(\gamma\); matched-line lifting would additionally prove the
transcendence of \(\gamma/\pi\), while the mixed route gives an all-order
linear-independence family.  No implication among the three inputs is known.
\end{remark}

\section{Conclusions and remaining problems}
\label{sec:conclusion}

The completed results fall into the following thematic groups.
\begin{enumerate}[label=\textup{(\arabic*)}]
 \item Every divisor \(\Ein-c\) has a universal polynomial inverse-trace
 calculus.  Its full trace ring is exactly \(F[c^{-1}]\), generated by the
 two even traces \(T_2,T_4\), while every transversal moving selection
 preserves algebraic independence.  The only nontrivial collision of
 \((T_2,T_3)\) is the pair \(\{6-2\sqrt3,6+2\sqrt3\}\); at every order
 there are also algebraic levels with zero inverse trace and wholly
 transcendental spectral divisors.  Finite-rank multiplicative groups have
 defect at most their rank and admit cofinite jointly independent
 subfamilies; the complete gamma-free, log-free cancellation space is
 identified.  The tails \(E_1(\alpha)\), and jointly the four classical
 families \(E_1,\Ei,\Ci,\operatorname{Si}\), have defect at most one more
 than the multiplicative rank.  Regularized lower incomplete-Gamma jets
 have defect at most the rank, uniformly across all jet orders, while fixed
 upper-jet rows have one further possible defect.  Sine-integral graph
 ideals, Dickman log-Laplace relations, and characteristic-zero algebraic
 matroids give further realizations of this engine.  At every finite set of
 levels, the complete trace field is exactly the rational quotient by the
 one-dimensional gauge shift.
 \item The dyadic and fixed-power balanced levels approach \(\gamma\) through
 explicit algebraically independent finite-stage data.  Their order-two
 and order-four traces give positive algebraically independent spectral
 decompositions of \(S_2\) and \(S_4\); the two positive carrier traces
 recover \(\gamma\).  The order-two carrier also gives an order-zero
 Laguerre--P\'olya determinant.  Neither the dyadic level sequence nor its
 second trace is P-recursive.
 \item The gamma-free map \(\mathscr R\) has transcendental inverse images at
 every algebraic level in its range, giving the explicit family \(r_m\).
 Normalized finite \(\Ein\)-certificates are unique.  This turns all
 cross-index real Euler witnesses into one countable exclusion theorem and
 yields the relative five-class theorem.  The critical cancellation
 coefficients \(u_N(q)\) have algebraic defect at most one.
 \item For the 2026 Challenge recurrence, the terminating \({}_2F_2\)
 state is explicitly gauge-equivalent to the contragredient normalized
 official Meijer--\(G\) state.  The finite bilinear identity fixes the
 conserved pairing; a flat dual-number \(SL_3\), equivalently \(SL_6\),
 includes the numerator.  The Beta--log transform connects the Challenge to
 the cubic family and every fixed offset.  A single rational carrier yields
 the denominator, numerator, full zeta log-jets, and a monic contact
 polynomial.  The exact maxima mixture, Meijer stop-loss interpretation,
 Poisson--binomial law, zero-density limit, regularized Gamma determinant,
 strengthened divisor, factorial-six Casoratian, all-initial-value limit
 classification, and integer solution lattice complete the finite
 structure.
 \item The Euler--Pad\'e convergents have a finite residue criterion on the
 \(p\)-adic unit boundary and converge as a full sequence at
 \(t=-1\) for \(p=2,3,5,13\).  One lcm-indexed rational subsequence converges
 placewise at the real place and at every fixed prime.  The denominator
 interpolation \(q_p\) identifies unbounded valuations with a \(p\)-adic
 zero, without confusing that event with failure of convergence.
 \item The dyadic augmentation tail is the unique minimum in the
 cumulative-positive cone, while outside that cone three points already
 produce dense tail values.  For fixed \(q\), the exact denominator
 transition for consecutive augmentation-one triples occurs at
 \(e^{(q-1)q^N}/N\).  A moving-point lower bound covering the displayed
 fixed system cannot have sublinear point cost; at linear cost the example
 forces the stated height loss.  The two-base form
 \(Y_N(2)-Y_{M_N}(3)\) retains the same leading \(2^N\) logarithmic scale.
 The exact asymmetric Sondow formula
 replaces the invalid universal entropy-floor heuristic.
 \item In each nondegenerate fixed residual direction of the Pr\'evost grid,
 the real Hilbert quotient is exactly saturated.  For the uncombined
 optimally saddled form, and absent additional exponential divisibility or
 cancellation, the baseline fixed-place ledger leaves the deficit
 \(2-\log(3+2\sqrt2)\).  The bounded-degree reduction transfers a surviving
 one-prime problem to quantitative approximation of \(\log2\), but does not
 close it.
 The fixed-prime Rodrigues deformation produces integral polynomials,
 exact logarithmic linear forms, positivity, and the endpoint factor
 \(p^{(1-\theta)n}\).  Even after the whole factor is favorably credited as
 fixed-place gain, the universal saddle-cancellation inequality shows that
 this ledger does not improve the Pr\'evost baseline.  Balanced factorial
 ratios themselves have only \(O(\log n)\) fixed-prime valuation.
 \item The displayed unframed lateral transition of the elementary Euler
 differential operator is gamma-free.  The missing information is the
 canonical framing.  In the mixed arithmetic-Gevrey route, the intersection
 conjecture \(\mathbf E\cap\mathbf D=\Qbar\) yields simultaneous linear
 independence of the dilation family and of all derivatives
 \(1,\Gamma'(1),\Gamma''(1),\ldots\).  These conclusions remain
 conditional.  In the Euler--Mascheroni motive the same extension
 coordinate is the off-diagonal entry of the period matrix.  The explicit
 matched-line lifting problem \cref{prob:euler-tate-fullness} would force
 the \(\Qbar\)-linear independence of \(1,\pi,\gamma\), not merely the
 transcendence of \(\gamma\).
\end{enumerate}

The principal quantitative unresolved problem can now be stated without
spectral language.

\begin{problem}[Common-coefficient separation]
\label{prob:common-coefficient}
Prove that for every \(s\in\Qbar\),
\[
 \limsup_{N\to\infty}
 \frac{-\log|2sY_N^2+Y_N-2|}{2^N}<1,
\]
or establish any equivalent separation that exploits the same coefficient
\(s\) across all \(N\).
\end{problem}

By \cref{thm:critical-slope}, this problem is equivalent to the
transcendence of \(S_2\), hence to the transcendence of \(\gamma\).  It is
therefore a precise quantitative target, not a smaller lemma already within
current fixed-point \(E\)-function theory.  Independently,
\cref{prob:euler-tate-fullness} is a categorical target for the canonical
rank-two Euler extension; its \(\Qbar\)-version would additionally prove
the transcendence of \(\gamma/\pi\).  The mixed-Gevrey intersection
conjecture gives a logically independent conditional route and the stronger
all-order linear-independence package, but it is not presently known.

We emphasize the corresponding negative conclusions.  The algebraic independence of the
finite stages does not determine the arithmetic nature of their limit.  The
Laguerre--P\'olya determinant is an exact spectral carrier but is not proved
arithmetic or \(D\)-finite; the regularized Gamma determinant does not prove
the irrationality of any odd zeta value.  Certificate uniqueness shows that producing more
linear witnesses of the same normalized type creates mutually exclusive
alternatives rather than independent chances to prove transcendence.  The
specific Pr\'evost and Rodrigues families analyzed here do not pay their
complete height ledger, even under the favorable fixed-place accounting
stated above.  The \(p\)-adic convergence theorems prove neither the
irrationality of \(E_p(-1)\), nor Kurepa's conjecture, nor full convergence
for every prime.  The external relation classifications do not close a
famous transcendence conjecture.  Finally, unframed Stokes data forget the
Euler extension coordinate, and the mixed-Gevrey theorems that recover it
are explicitly conditional.
Accordingly, no unconditional theorem in this paper proves the irrationality or
transcendence of \(\gamma\), \(\delta\), \(S_2\), or any odd zeta value; the
exact Euler apertures above state what is still missing.
